
% 03. The Illusion of Expansion and the Reinterpretation of the Big Bang.tex

\section{팽창의 환상과 빅뱅의 재해석}
우리가 시간과 공간이라고 부르는 것이 우주의 심연을 가리는 생존용 '가짜 바탕화면'에 불과하다면, 그 얄팍한 화면을 걷어냈을 때 텅 빈 무대 위에 드러나는 우주의 실제 하드웨어는 과연 무엇인가.
디스턴스(DISTANCE) 에세이는 그 근본에 '에너지 갭(Energy Gap)'이 존재한다고 장엄하게 선언한다. 우주에서 발생하는 모든 운동과 변화, 그리고 사우리는 빅뱅과 우주의 팽창을 이야기할 때 흔히 '부풀어 오르는 풍선'의 비유를 배운다. 작은 풍선 표면에 찍혀 있던 점들이, 풍선이 부풀어 오름에 따라 서로 점점 멀어지는 현상. 현대 물리학은 이를 근거로 우주를 채우고 있는 '공간 자체가 기하학적으로 늘어나고 있다'고 설명하며, 우리는 그 장엄한 공간 팽창의 서사를 한 치의 의심 없이 절대적인 우주적 진리로 받아들여 왔다.

그러나 디스턴스(DISTANCE)의 렌즈를 들이대면, 이 지배적인 서사 안에는 치명적이고도 거대한 가정이 숨어 있음이 폭로된다. 그것은 '공간 자체가 실제로 늘어났다'는 가정이다. 과연 그러한가? 우리는 단 한 번도 공간이라는 투명한 그릇 자체가 고무줄처럼 늘어나는 것을 실제로 관측한 적이 없다. 인간이 천체망원경을 통해 실제로 관측하는 물리적 팩트는 오직 은하 사이의 '평균 밀도 변화', '에너지의 분포', 그리고 '물질 농도의 희박해짐'뿐이다. 우리가 관측하는 모든 것은 결국 구조와 밀도의 역동적인 변화일 뿐인데, 시공간을 절대적 하드웨어로 맹신하는 인간의 뇌가 이를 '공간 자체의 변형(팽창)'으로 착각하여 번역해 내고 있는 것이다.

이 기만적인 착각을 완벽하게 부수기 위해, 투명한 물이 가득 담긴 컵의 비유를 들어 보자. 컵의 중앙에 검은 잉크 한 방울을 떨어뜨린다. 아주 짧은 찰나, 극히 작은 영역에 엄청난 농도의 색소가 뭉쳐 있는 초고밀도 상태가 형성된다
. 디스턴스의 역학적 언어로 번역하면, 잉크가 뭉쳐 있는 곳과 주변의 맹물 사이에 극단적인 '거대한 에너지 갭(농도 차이)'이 존재하며, 이를 해소하려는 맹렬한 '모멘텀'이 끓어오르는 상태이다.
시간이 지나면 잉크는 필연적으로 주변으로 확산되며 뻗어 나간다. 더 넓은 영역으로 흩어지며 고밀도 상태는 무너지고, 결국 컵 전체가 균일한 회색의 저밀도 상태로 뒤섞인다. 바로 이 지점에서 가장 예리한 질문을 던져야 한다. 잉크가 퍼지는 동안 컵(공간)이 커졌는가? 결코 아니다. 변한 것은 컵의 크기가 아니라 오직 '밀도의 분포'뿐이다. 그러나 만약 컵이라는 외부 경계를 인지하지 못하고 오직 잉크 입자들 시점에서만 세상을 관찰하는 내부의 관찰자가 있다면, 그는 사방으로 멀어지는 입자들을 보며 확신에 차 이렇게 선언할 것이다. "모든 입자가 서로 멀어지고 있다. 즉, 공간이 팽창하고 있다!"

이 완벽한 논리적 오류가 바로 오늘날 인류가 빅뱅과 우주 팽창을 해석하는 방식이다. 컵의 크기가 1리터이든, 태평양만 하든, 심지어 경계가 존재하지 않는 '무한대'이든 현상의 본질은 단 한 치도 달라지지 않는다. 잉크는 무한한 컵 속에서도 끊임없이 퍼져나가고 농도는 희박해지며, 맹렬한 평활화는 멈추지 않는다. 크기가 유한하든 무한하든 관측자는 여전히 입자들이 멀어지는 환상을 보게 되며, 본질은 '용기의 팽창'이 아니라 오직 '모멘텀의 해소'라는 사실이다.
이 서늘한 역학적 진실 앞에서 '빅뱅(Big Bang)'의 의미는 완전히 재구성된다. 

빅뱅 이전의 우주는 상상조차 할 수 없는 초고밀도로 모든 것이 응축된 궁극의 에너지 아일랜드였다. 그 상태는 우주 역사상 가장 거대한 에너지 갭이자 가장 폭발적인 모멘텀을 품고 있는 상태였다. 그렇게 극단적인 모멘텀은 결코 안정적으로 머물 수 없으며, 열역학적 필연에 따라 자신의 격차를 해소하려는 맹렬한 분산(평활화)을 시작할 수밖에 없다.
우리가 '빅뱅'이라 부르는 현상은 빈 공간 속에서 물질이 기하학적으로 폭발하며 공간을 찢어 늘린 사건이 아니다. 그것은 초고밀도의 에너지 아일랜드가 스스로 품고 있던 극단적인 모멘텀을 무너뜨리며, 궁극의 평활화를 향해 우주적 잉크를 흩뿌리기 시작한 '거대한 모멘텀 해소의 출발점'인 것이다.

현대 우주론이 공간 팽창의 가장 강력한 증거로 제시하는 '적색편이(Red Shift)' 현상 역시 마찬가지다. 디스턴스의 렌즈로 보면, 적색편이는 기하학적인 공간이 늘어나 파장이 길어진 마법이 아니다. 그것은 에너지가 평활화를 향해 무너져 내리며 분포가 점점 희박해지는 과정, 즉 존재들 사이의 '역학적 디스턴스(관계적 거리)'가 맹렬하게 증가하고 모멘텀이 소진되어 가는 과정에서 나타나는 열역학적 결과물일 뿐이다.
결국 우주의 역사는 텅 빈 3차원 공간이 풍선처럼 커져 온 기하학의 역사가 아니다. 그것은 초고도의 밀도가 끝을 알 수 없는 옅음(평활화)을 향해 맹렬하게 흩어지며 모멘텀을 소진해 온 '추락과 분산의 역사'이다. 오늘날 우리가 경이롭게 바라보는 은하와 별, 행성, 그리고 생명과 문명조차 이 거대한 잉크의 확산(평활화) 과정 속에서 극히 일시적으로 뭉쳐진 국소적인 에너지 아일랜드(소용돌이)들에 불과하다. 언젠가 잉크가 물과 완벽히 섞여 흐름을 멈추듯, 이들 또한 궁극적으로는 더 큰 스케일의 완벽한 평활화 속으로 남김없이 흡수될 것이다.

우주가 커지고 있기 때문에 별들이 멀어지는 것이 아니다. 우주가 완벽한 정지(엔트로피 최대화)를 향해 스스로를 맹렬하게 '섞고' 있기 때문에, 존재들 간의 거리가 멀어지는 것이다. 우주는 자신의 거대한 모멘텀을 쉴 새 없이 해소하며, 단 1초의 멈춤도 없이 거대한 평활화의 폭포수 아래로 스스로를 내던지고 있다.
